% Image credits for the photographs and AI illustrations of Book 3
% (University Physics, Year 1). Machine-readable license records live in
% images/book3/CREDITS.md; keep the two files in sync.
\chapter*{Beeldverantwoording}

De tekeningen in dit boek zijn origineel. De foto's worden gebruikt
onder hun respectieve vrije licenties, met dank aan hun makers:

\begin{itemize}
  \item \emph{Kibblebalans (NIST-4)} (hoofdstuk 1): J.~L.~Lee voor
        NIST, publiek domein.
  \item \emph{Yerkes-refractor van 40 inch} (hoofdstuk 4): Kb9vrg,
        Wikimedia Commons, publiek domein.
  \item \emph{De Maan, naar de Aarde gekeerde zijde} (hoofdstuk 4):
        NASA/GSFC/Arizona State University (Lunar Reconnaissance
        Orbiter), publiek domein.
  \item \emph{Gitaarsnaren} (hoofdstuk 5): Acagastya, Wikimedia
        Commons, CC0.
  \item \emph{Oscilloscoop, sinus- en blokgolf} (hoofdstuk 8): Xato,
        Wikimedia Commons, publiek domein.
  \item \emph{Looping van een achtbaan} (hoofdstuk 11): Jeremy
        Thompson, Wikimedia Commons, CC~BY~2.0.
  \item \emph{Internationaal ruimtestation} (hoofdstuk 12): NASA,
        publiek domein.
  \item \emph{GPS-satelliet} (hoofdstuk 16): NASA, publiek domein.
  \item \emph{Cyclotron van 60 inch in Berkeley} (hoofdstuk 17):
        Amerikaanse ministerie van Energie, publiek domein.
  \item \emph{Slinger van Foucault, Panthéon} (hoofdstuk 18): Olga
        Khomitsevich, Wikimedia Commons, CC~BY~2.0.
  \item \emph{Onderzeeboot die opduikt} (hoofdstuk 21): US Navy,
        publiek domein.
  \item \emph{Bliksem} (hoofdstuk 26): Mark Coldren, Wikimedia
        Commons, CC0.
  \item \emph{Millikans oliedruppelopstelling} (hoofdstuk 27):
        Daderot, Wikimedia Commons, CC0.
  \item \emph{Kompas} (hoofdstuk 28): Evan-Amos, Wikimedia Commons,
        publiek domein.
  \item \emph{Hoogspanningslijn} (hoofdstuk 28): Stefan Andrej
        Shambora, Wikimedia Commons, CC~BY~2.0.
  \item \emph{MRI-scanner} (hoofdstuk 29): Jan Ainali, Wikimedia
        Commons, CC~BY~3.0.
  \item \emph{Emissiespectrum van waterstof} (hoofdstuk 30):
        Merikanto en Adrignola, Wikimedia Commons, CC0.
\end{itemize}

Volledige bronvermeldingen en licentieadressen van elke foto staan in
\texttt{images/book3/CREDITS.md} in de repository van het boek.

\bigskip
Naast de foto's zijn vierentwintig figuren in de hoofdstukken 1--4,
6--9, 12--15, 19, 21--25 en 27--30 door kunstmatige intelligentie
gemaakte illustraties, vervaardigd met de beeldgeneratie van OpenAI op
basis van prompts die de redacteuren van het boek hebben geschreven,
en vóór opname nagelezen op natuurkundige juistheid. De volledige
prompt van elk gegenereerd beeld staat in
\texttt{images/book3/ai/PROMPTS.md} in de repository.
\clearpage
